\section{Discussion}
\label{sec:discussion}

\begin{table}[t]
    \centering
    \caption{Verdict for each part of the hypothesis.}
    \label{tab:verdict}
    \renewcommand{\arraystretch}{1.2}
    \resizebox{\textwidth}{!}{%
    \begin{tabular}{@{}llp{9.2cm}@{}}
        \toprule
        \textbf{Claim} & \textbf{Verdict} & \textbf{Evidence} \\
        \midrule
        H1a: gains are logarithmic & Partly (math only) & Log is the best model for math but explains little variance ($R^2 = 0.28$). Code follows a bounded saturating curve; creative ``gains'' are negative. \\
        H1b: 80\% of benefit at \ratio{3} & Not supported & Code 63\% (CI 13--100\%), math 40\% (CI 6--64\%), undefined for creative writing. \\
        H2a: plateau or degradation beyond \ratio{10} & Depends on domain and metric & Code accuracy plateaus; math accuracy shows no plateau up to \ratio{50}; creative fit degrades. NLL on real references degrades in all three. \\
        H2b: mode collapse is the mechanism & Consistent, not proven & Pool diversity saturates and students inherit teacher modes, but we do not manipulate diversity. \\
        H3: code most resilient, creative earliest & Half supported & Creative writing degrades from the first synthetic example. Code is stable, but math tolerates and benefits from the highest ratios. \\
        \bottomrule
    \end{tabular}
    }
\end{table}

\Tabref{tab:verdict} summarizes the verdicts.
We now interpret the pattern and state what the study cannot show.

\para{Two effects pull in opposite directions.}
Synthetic data transfers the teacher's \emph{skill}, which helps where an answer is right or wrong, and the teacher's \emph{distribution}, which hurts where the target is the human distribution.
Math shows both at once: the student becomes a better solver while it becomes a worse imitator of the reference style of \gsm.
Stories have no skill component in our metrics, so only the distribution shift is visible.
This view explains why the task metric and the shared metric disagree beyond \ratio{10}, and why no single ceiling exists.

\para{Why does synthetic math data beat real data?}
At the \ratio{30} size the synthetic mix is 6.5 points more accurate than 3,100 real \gsm examples.
We did not expect this.
A plausible explanation, which we did not test, is that the teacher's worked solutions are more uniform and explicit than human-written ones, and that a 0.5B model learns from them more easily.
We also did not verify the correctness of the teacher's math solutions.

\para{Why does code saturate early?}
The code pools are the most collapsed, with 79\% near-duplicates at 3,000 examples, so ratios beyond about \ratio{5} add almost no new problems.
The student keeps re-reading \texttt{factorial}.
The execution filter may contribute by favoring easy problems the teacher can solve, but our ablation cannot settle this.

\para{Small doses can be the worst doses.}
On the shared metric, \ratio{1} is the worst ratio for code and is worse than every ratio from \ratio{2} to \ratio{10} for stories.
One reading, which we did not test, is that a small amount of synthetic data shifts the distribution without yet adding much skill.

\para{Diversity metrics can miss homogenization.}
Distinct-2 of student stories does not fall with the synthetic ratio, while the Fr\'echet distance to human stories roughly doubles and one teacher-preferred name becomes 15 times more frequent.
Studies that track only $n$-gram diversity would miss this.

\para{Practical implications.}
For verifiable skills such as code and arithmetic reasoning, extra synthetic data did not hurt task accuracy up to \ratio{30}--\ratio{50} in our setting, but most of the code benefit arrived by about \ratio{5}.
For open-ended generation, where the goal is to match a human distribution, synthetic data from a single generator hurt at every ratio we tested, and 100 more real examples were worth more than 5,000 synthetic ones.

\subsection{Limitations}
\label{sec:limitations}

\para{Statistical power.} Three seeds and 500 test items give 95\% CIs of about $\pm 4$ points on accuracy. Differences of under about 4 points between adjacent ratios are within noise. Only the \ratio{3}-against-0 and peak-against-0 contrasts are significant for code and math, so ``plateau'' for code and ``still rising'' for math are statements our data cannot sharply distinguish. A seed-level bootstrap with three seeds is coarse. The NLL-based results are much tighter.

\para{The share statistic.} $B(3)$ uses the maximum over a noisy grid, which biases the share at \ratio{3} downward. The location of the code maximum is itself uncertain: \ratio{30} is the best ratio in only 39\% of bootstrap resamples. For math the maximum is at the edge of the grid.

\para{One configuration.} We use one teacher (3B), one student (0.5B, LoRA rank 32), one real-set size (100), and one generation recipe (3-shot, $T = 1$, no de-duplication, no diversity prompting). A stronger teacher, attribute-controlled prompting~\citep{yu2023attributed}, or de-duplication would likely move the curves.

\para{Training budget.} We train for at most two epochs with early stopping, so larger ratios receive more gradient steps. We did not match compute across ratios.

\para{Code is not directly comparable to the other domains.} Its pool is execution-filtered, its grid stops at \ratio{30}, and its real-data control stops at \ratio{2.74} because \mbpp is small.

\para{Creative quality.} We measure fit to human text and distributional distance. We collected no human or model-judge ratings of story quality; synthetic-trained stories may read as more polished while being less human-like.

\para{Contamination.} \mbpp, \gsm, and \roc are probably in the pre-training data of both models. This affects absolute levels more than the comparison between ratios.

\para{Zero-shot math baseline.} The 256-token generation cap depresses the zero-shot math score; a pilot with 320 tokens on all 1,319 test items gave 41.8\%.

\para{Mechanism.} The link between mode collapse and the dose--response curves is correlational. No experiment varied pool diversity while holding size fixed.
